\section{Analytic Tools}
\label{sec:core}

Once a corpus has produced an accepted fitted chain, the question is
what that chain answers without another benchmark run. Three
absorbing-chain tools address the deployment questions of
\S\ref{sec:intro}: reliability at a horizon, sensitivity to a local
change, and compatibility among common benchmark metrics.

\subsection{Closed-Form Reliability}

The first deployment question is reliability under a step budget:
starting at $s_0$, what is the chance of reaching $\oplus$ by step
$d$? In the running loop, this asks whether the agent reaches a
correct final answer within $d$ reasoning/tool/observation steps.

\begin{proposition}[Closed-form reliability; Kemeny--Snell~\cite{kemenysnell}]
\label{thm:closed}
Under Assumptions~\ref{A:trans}--\ref{A:init}, for every
$d\in\mathbb{N}_{\ge 0}$,
\begin{equation}
R(d) = e_{s_0}^\top(I-Q^{d})\,N\,R_\oplus,
\qquad N=(I-Q)^{-1},
\label{eq:t1finite}
\end{equation}
and
$R_\infty = \lim_{d\to\infty}R(d) = e_{s_0}^\top N R_\oplus$.
\end{proposition}

\noindent\emph{Proof idea.}
Summing all paths that spend $t=0,\ldots,d-1$ steps in transient states
and then exit to success gives
$R(d)=e_{s_0}^{\top}\sum_{t=0}^{d-1}Q^tR_\oplus$; the identity
$\sum_{t=0}^{d-1}Q^t=(I-Q^d)(I-Q)^{-1}$ yields \eqref{eq:t1finite}, and
transience gives $Q^d\to0$. Since $N_{ij}$ is the expected number of
visits to state $j$ before absorption, the closed form separates two
effects an operator can inspect: how often the agent visits each state,
and how likely success is after those visits.

\subsection{Perturbation Sensitivity}

The second question is counterfactual: if a local transition changes,
how much can end-to-end reliability move? A new fallback tool may
change one row of $Q$ by redirecting failed tool calls back to
planning, and the fitted chain gives a first-order sensitivity
calculation before any benchmark rerun.

\begin{proposition}[Perturbation sensitivity; Stewart~\cite{stewart1998}]
\label{thm:perturb}
Let $Q(\varepsilon)=Q_0+\varepsilon\Delta$ be the perturbation family
from Definition~\ref{def:perturb}. Write
$N_0=(I-Q_0)^{-1}$ and let $R_{\oplus,0}$ be the fixed success-exit
vector. Assume
$\varepsilon_{\max}\|N_0\|_\infty\|\Delta\|_\infty<1$.
For every $0\le\varepsilon\le\varepsilon_{\max}$,
\begin{equation}
\bigl|R_\infty(\varepsilon)-R_\infty(0)\bigr|
\le
\varepsilon \|N_0\|_\infty^2\|\Delta\|_\infty
        \|R_{\oplus,0}\|_\infty
 + C\varepsilon^2,
\label{eq:t2}
\end{equation}
with
$C=\|N_0\|_\infty^3\|\Delta\|_\infty^2\|R_{\oplus,0}\|_\infty
 /(1-\varepsilon_{\max}\|N_0\|_\infty\|\Delta\|_\infty)$. Moreover,
\begin{equation}
R_\infty(\varepsilon)-R_\infty(0)
=\varepsilon e_{s_0}^{\top}N_0\Delta N_0R_{\oplus,0}
 + O(\varepsilon^2).
\label{eq:t2lin}
\end{equation}
\end{proposition}

\noindent\emph{Proof idea.}
Expand $(I-Q(\varepsilon))^{-1}$ with the Neumann series around $Q_0$;
the smallness condition bounds the remainder by the displayed constant.
In the first-order term $N_0\Delta N_0$, the left factor counts how
often execution reaches the changed row and the right factor measures
success value after the perturbed transition, so a large
$\|N_0\|_\infty$ warns that small local changes can have large global
effects. The bound is best read as a screening calculation that
identifies local changes meriting a new benchmark run. As stated it
applies to perturbations of $Q$ with $R_\oplus$ fixed; changes to
success exits require the corresponding linear term for the exit
vector, and the real-data fallback-tool result of \S\ref{sec:ss14} is
of that second kind, so we evaluate it exactly from
Proposition~\ref{thm:closed} rather than through this bound.

\subsection{Metric Unification}
\label{sec:unify}

The third question is semantic. pass$@k$, pass$^k$, and the RDC often
appear to be competing summaries, but under the fitted chain they are
different views of one first-passage distribution: the disagreement is
about which projection to report, not about the underlying reliability
event. We keep the benchmark name RDC, but the quantity modeled here is
the cumulative success probability by horizon $d$.

\begin{theorem}[Metric unification]
\label{thm:unify}
Let $R_\infty=e_{s_0}^{\top}NR_\oplus$. Under the i.i.d.\
replication assumption A\ref{A:iid},
\begin{align}
\mathrm{pass}^{k} &= R_\infty^{k}, \label{eq:t3-passk}\\
\mathrm{pass}@k  &= 1-(1-R_\infty)^{k}, \label{eq:t3-passatk}\\
\mathrm{RDC}(d)  &= R(d)
 = e_{s_0}^{\top}(I-Q^{d})NR_\oplus. \label{eq:t3-rdc}
\end{align}
Thus the three metrics are restrictions or transformations of the
same first-passage distribution.
\end{theorem}

\noindent\emph{Proof idea.}
pass$^k$ is the probability that all $k$ independent trials succeed,
pass$@k$ is one minus the probability that all $k$ fail, and the RDC is
the finite-horizon reliability of Proposition~\ref{thm:closed}.

Repeated agent runs can, however, share latent conditions such as a
common prompt template, cache state, or task difficulty. The next
result makes that caveat measurable: it shows how shared variation
changes the two repeated-trial metrics even when marginal reliability
remains $R_\infty$.

\begin{theorem}[Correlated-trial inequalities]
\label{thm:corr}
Suppose the $k$ trials share a latent state $\xi$ with
$\Pr[X_i=1\mid\xi]=p(\xi)$ and
$\mathbb{E}_\xi[p(\xi)]=R_\infty$, and are independent conditional on
$\xi$. Then
\begin{align}
\mathrm{pass}^{k} &= \mathbb{E}_\xi[p(\xi)^k]
                  \ge R_\infty^k, \label{eq:t3p-passk}\\
\mathrm{pass}@k  &= 1-\mathbb{E}_\xi[(1-p(\xi))^k]
                  \le 1-(1-R_\infty)^k. \label{eq:t3p-passatk}
\end{align}
For $k\ge2$, equality holds iff $p(\xi)$ is almost surely constant.
\end{theorem}

\noindent\emph{Proof idea.}
Condition on the latent variable and apply Jensen's inequality to
the convex functions $p^k$ and $(1-p)^k$.

\begin{corollary}[Diagnostic]
\label{cor:diag}
The gap
$\widehat{\mathrm{pass}^{k}}-\hat R_\infty^{\,k}$ is a
lightweight diagnostic for hidden cross-trial correlation. Large
positive gaps indicate that repeated trials are not behaving as
i.i.d.\ samples from one fixed chain.
\end{corollary}

Hidden heterogeneity thus raises all-success estimates and lowers
at-least-one-success estimates relative to the i.i.d.\ formulas, so a
large gap should trigger inspection before repeated samples are treated
as independent trials: an assumption behind pass$@k$ and pass$^k$
becomes a diagnostic rather than a hidden convention.
